\section{Contextual Bandits}
% Slide 34
\begin{frame}[shrink=10]{Contextual Bandits}
\begin{itemize}
    \item A Contextual Bandit is a 3-tuple $(\mathcal{A}, \mathcal{S}, \mathcal{R})$
    \item $\mathcal{A}$ is a known set of $m$ actions (``arms'')
    \item $\mathcal{S} = \mathbb{P}[s]$ is an unknown distribution over states (``contexts'')
    \item $\mathcal{R}^{a}_{s}(r) = \mathbb{P}[r \mid s, a]$ is an unknown probability distribution over
          rewards
    \item At each step $t$, the following sequence of events occur:
    \begin{itemize}
        \item The environment generates a states $s_t \sim \mathcal{S}$
        \item Then the AI Agent (algorithm) selects an actions $a_t \in \mathcal{A}$
        \item Then the environment generates a reward $r_t \in \mathcal{R}^{a_t}_{s_t}$
    \end{itemize}
    \item The AI agent's goal is to maximize the Cumulative Reward:
    \[
        \sum_{t=1}^{T} r_t
    \]
    \item Extend Bandit Algorithms to Action-Value $Q(s, a)$ (instead of $Q(a)$)
\end{itemize}
\end{frame}
